\chapter{Suspicious Website Logins}
\label{ch:1.2}
\chaptertopic{Topic 1.2}

\chapterquote{If you think technology can solve your security problems, then you don't understand the problems and you don't understand the technology.}{Bruce Schneier, \emph{Secrets and Lies} (2000)}

\section{Motivation}

Chapter~\ref{ch:1.1} was about tricking a human into handing over a credential. This chapter is
about what happens when an attacker skips the human entirely and attacks the login form itself.
Just as important is how a defender reads the evidence that attack leaves behind. A
password is not just a secret; it is a claim about probability. Every password you choose sits
somewhere on a spectrum from ``a computer can guess this in under a second'' to ``the sun will
burn out before a computer guesses this by chance.'' This chapter builds the vocabulary to
locate a password on that spectrum, to recognize an attack against one in progress, and to
understand why a second factor changes the entire calculation.

\begin{apalignment}{Unit 1: Introduction to Security}{Topic 1.2 -- Suspicious Website Logins}
\begin{apobjectives}
\apobjective{1.2.A}{Analyze Risk}{Recognize the common signs that an account is the target of a password attack.}
\apobjective{1.2.B}{Analyze Risk}{Explain how adversaries take advantage of weak authentication practices.}
\apobjective{1.2.C}{Mitigate Risk}{Explain the measures that make authentication stronger.}
\end{apobjectives}
\end{apalignment}

\section{Recognizing an Attack in the Logs}

An automated login attack does not look like a human typing a password wrong once. It looks like
a machine trying hundreds of guesses per minute, and the traces it leaves in an authentication
log are distinctive once you know what to look for.

\begin{definition}{Brute-force attack}
An automated, exhaustive attempt to guess a credential by systematically trying many candidate
values, typically far faster than a human could type them.
\end{definition}

Three signals recur across real intrusion-detection systems, and they are far more reliable
together than any one is alone.

\textbf{Repeated failed attempts.} A single failed login is background noise -- everyone
mistypes a password occasionally. Five failures for the same account inside sixty seconds is a
widely used threshold instead. It is fast enough to catch automated tools without flooding a
security team with false positives from ordinary human error.

\textbf{Unusual times.} People behave in habitual patterns. An account that normally shows
activity only between 08:00 and 18:00 on weekdays, and then shows a successful login at 02:47 on
a Saturday, has either an unusually dedicated user or a compromised account. Systems that build a
per-account behavioral baseline can flag this deviation automatically.

\textbf{Unknown devices.} Most authentication systems track which devices (by MAC address, IP,
or device fingerprint) have successfully logged in historically. A brand-new device succeeding
immediately after a burst of failures is close to a confession.

\begin{example}[Why correlation beats any single signal]
Each signal alone is weak evidence. A new device might just be a new phone. An odd login time
might be a night-shift worker. A few failures might be a typo. But an account showing five
failures at 3~a.m., from a device that has never appeared before, followed immediately by a
success, is a pattern that an innocent explanation struggles to account for. Combining
independent weak signals into one strong verdict is called \textbf{correlation}. It is the
organizing principle behind nearly every detection rule you will meet for the rest of this
course -- physical, network, or application layer.
\end{example}

\begin{lstlisting}[language=Python, caption={Flagging accounts under brute-force attack from a stream of login events}]
from collections import defaultdict

def flag_bruteforce(events, window_seconds=60, threshold=5):
    """events: list of dicts like
       {"user": str, "timestamp": int, "result": "FAILED" or "SUCCESS"}
       Returns the sorted set of usernames with >= threshold failures
       inside any window_seconds-wide sliding window."""
    by_user = defaultdict(list)
    for e in events:
        if e["result"] == "FAILED":
            by_user[e["user"]].append(e["timestamp"])

    flagged = set()
    for user, times in by_user.items():
        times.sort()
        left = 0
        for right in range(len(times)):
            while times[right] - times[left] > window_seconds:
                left += 1
            if right - left + 1 >= threshold:
                flagged.add(user)
                break
    return sorted(flagged)
\end{lstlisting}

\section{Evaluating Password Strength Like an Attacker}

Detecting an attack in progress is a defensive skill. Understanding why a given password falls
quickly is a design skill, and it requires seeing the login form the way the attacker does.

\begin{definition}{Dictionary attack}
A brute-force strategy that replaces exhaustive random guessing with structured guessing from a
wordlist of likely candidate passwords, ordered roughly by real-world frequency.
\end{definition}

Attackers rarely start from \texttt{aaaaaaaa} and count upward -- the search space is too large.
Instead they start from what has already worked before: leaked breach data, target-specific
information harvested from social media, and \textbf{mangling rules}. The 2009 RockYou leak alone
contributed fourteen million real passwords to modern wordlists. Mangling rules mechanically
transform every dictionary word into hundreds of plausible variants: capitalize the first letter,
append \texttt{1} or \texttt{!}, substitute \texttt{a} with \texttt{@}. A password like
\texttt{P@ssw0rd!} feels custom-made to a human but is one of the very first variants any
mangling engine tries.

Four weaknesses are worth checking for explicitly, because they compound:

\begin{enumerate}
\item \textbf{Too short.} This course uses a 12-character floor. Each additional character
multiplies the total search space by a full factor of the alphabet size. Going from 8 to 12
characters, for instance, increases a mixed-character search space by roughly seven orders of magnitude.
\item \textbf{Dictionary membership.} A password that appears verbatim, or as a simple mangled
variant, in a known wordlist is cracked in the first pass regardless of any other property it has.
\item \textbf{Restricted character set.} A password built only from lowercase letters is drawn
from a pool of 26 symbols per position rather than 90-plus; attackers who suspect this narrow the
search accordingly.
\item \textbf{Sequential or repeated patterns.} \texttt{qwerty}, \texttt{aaa111}, and
\texttt{abcabc} are explicitly enumerated in every modern wordlist and guessed early.
\end{enumerate}

\begin{misconception}
\textbf{``Adding a number and a symbol makes a password secure.''} If the underlying word is
already in the dictionary, mangling rules generate the ``clever'' variant automatically --
attackers see \texttt{p@ssword} as a five-second guess, not a hard problem. Unpredictability and
length are what generate real entropy; a predictable substitution pattern generates almost none.
\end{misconception}

\section{Why Length Beats Complexity}

The math behind the previous section's advice is worth making explicit. A password drawn
uniformly at random from an alphabet of size $A$, with length $L$, has a search space of $A^{L}$
possible values. Widening the alphabet from, say, 26 lowercase letters to 94 printable ASCII
characters multiplies the space by roughly a factor of $(94/26)$ for every position -- a modest,
one-time gain. Adding a single additional \emph{character} to the password, by contrast,
multiplies the entire space by another full factor of $A$. Length compounds; complexity merely
scales. This is precisely why NIST SP 800-63B moved away from mandatory complexity rules: they
push users toward predictable, mangling-rule-friendly patterns like \texttt{Password1!}. Such a
pattern offers a smaller entropy gain than four random, unrelated words strung together as a
passphrase.

\section{Multi-Factor Authentication and TOTP}

Even a genuinely long, random password is still just a single secret. If that secret is
phished, leaked in a breach, or observed over someone's shoulder, the attacker gains permanent
access with no further work required. \textbf{Multi-Factor Authentication (MFA)} closes this gap
by requiring proof of identity from two or more independent categories: something you
\emph{know} (a password), something you \emph{have} (a phone, a hardware token), or something
you \emph{are} (a fingerprint). An attacker who has only compromised one category still has
nothing.

\begin{definition}{Time-based One-Time Password (TOTP)}
A six-digit authentication code derived from \texttt{HMAC-SHA1(secret, time\_step)}, where
\texttt{time\_step} is the current 30-second interval since the Unix epoch. Defined in RFC~6238,
it is the algorithm behind most authenticator-app second factors.
\end{definition}

TOTP's defining property is that it requires \emph{no network communication at verification
time}. The server and the authenticator app each independently hold the same shared secret $K$,
established once during enrollment (typically transmitted over HTTPS as a QR code). From that
point on, both sides compute
\[
T = \left\lfloor \frac{\text{unix\_timestamp}}{30} \right\rfloor, \qquad
\text{code} = \operatorname{Truncate}\big(\mathrm{HMAC\text{-}SHA1}(K,\, T)\big) \bmod 10^{6},
\]
and compare. The code changes every 30 seconds, and the space of possible codes is only
one million values. An intercepted code becomes useless within moments -- a fundamentally
different security posture from an SMS code, which travels across a carrier network that can be
intercepted or redirected via SIM swapping.

\begin{figure}[h]
\centering
\begin{tikzpicture}[
  box/.style={draw=kickstartermaroon, thick, rounded corners, minimum width=3.4cm, minimum height=1cm, align=center, font=\small},
  lbl/.style={font=\scriptsize\sffamily}
]
\node[box, fill=kickstartercream] (client) at (0,2.2) {Authenticator app\\shares secret $K$};
\node[box, fill=kickstartercream] (server) at (6,2.2) {Server\\shares secret $K$};
\node[box] (ct) at (0,0.9) {compute $T = \lfloor t/30 \rfloor$};
\node[box] (st) at (6,0.9) {compute $T = \lfloor t/30 \rfloor$};
\node[box] (cc) at (0,-0.4) {HMAC-SHA1$(K,T)$\\$\to$ 6-digit code};
\node[box] (sc) at (6,-0.4) {HMAC-SHA1$(K,T)$\\$\to$ 6-digit code};
\node[box, fill=kickstartermaroon!15] (cmp) at (3,-1.9) {codes equal?\\\textbf{accept / reject}};
\draw[-{Latex[length=2mm]}] (client) -- (ct);
\draw[-{Latex[length=2mm]}] (server) -- (st);
\draw[-{Latex[length=2mm]}] (ct) -- (cc);
\draw[-{Latex[length=2mm]}] (st) -- (sc);
\draw[-{Latex[length=2mm]}] (cc) -- (cmp);
\draw[-{Latex[length=2mm]}] (sc) -- (cmp);
\node[lbl, gray, align=center] at (3,2.2) {\emph{no network traffic between}\\\emph{these two computations}};
\end{tikzpicture}
\caption{Both sides of a TOTP exchange compute the code independently from the same shared secret
and the same 30-second time step -- only the final comparison, at login, ever crosses the network.}
\label{fig:1.2-totp-flow}
\end{figure}

\begin{lstlisting}[language=Python, caption={A from-scratch TOTP verifier using only the standard library}]
import hmac, hashlib, struct

def totp_code(secret_bytes: bytes, timestamp: int, step: int = 30, digits: int = 6) -> str:
    """Compute the TOTP code for the 30-second window containing `timestamp`."""
    counter = timestamp // step
    msg = struct.pack(">Q", counter)          # 8-byte big-endian time step
    digest = hmac.new(secret_bytes, msg, hashlib.sha1).digest()
    offset = digest[-1] & 0x0F                # dynamic truncation (RFC 4226)
    truncated = struct.unpack(">I", digest[offset:offset + 4])[0] & 0x7FFFFFFF
    return str(truncated % (10 ** digits)).zfill(digits)

def verify_totp(secret_bytes: bytes, submitted_code: str, timestamp: int) -> bool:
    return hmac.compare_digest(totp_code(secret_bytes, timestamp), submitted_code)
\end{lstlisting}

Once you have computed a TOTP code by hand from raw bytes, ``the authenticator app just knows the
code'' stops being a black box. It becomes an equation you can evaluate yourself -- which is
precisely why this course asks you to implement it rather than only use it.

\begin{misconception}
\textbf{``MFA is unbreakable.''} A phished TOTP code can still be relayed to the real login form
in real time by an attacker fast enough to beat the 30-second window. MFA raises the bar
dramatically -- an attacker now needs two independent compromises instead of one -- but it does
not eliminate social engineering; it changes what the attacker has to accomplish through it.
\end{misconception}

\section{Case Studies in Authentication Failure}

\begin{example}[Colonial Pipeline, 2021: One Dormant Account]
In May 2021 the operator of the largest refined-fuel pipeline in the United States was forced to
shut down 5,500 miles of pipeline after a ransomware deployment by the DarkSide group. The initial
entry point was a legacy VPN account: one that was not meant to still be in use, had never been
formally decommissioned, and was protected only by a single password with no multi-factor
authentication behind it at all.\srccite{1.2}{1} Mandiant's
investigation traced that password to an earlier, unrelated data breach -- the employee had almost
certainly reused it on a different, already-compromised service. The company paid a ransom of
roughly \$4.4~million in bitcoin, of which federal investigators later recovered about
\$2.3~million by seizing the attacker's wallet key.\srccite{1.2}{2} The lesson is not simply
``enable MFA'' -- true as that is. It is that an account nobody uses is also an account nobody is
watching. It appears in no active-user risk review, and its compromise triggers no behavioral
anomaly, because there is no established baseline of normal behavior to deviate from.
Decommissioning credentials on employee or system departure, and automatically disabling any
account idle beyond a defined threshold, closes a gap that technical monitoring alone cannot.
\end{example}

\begin{example}[Cloudflare vs.\ Twilio, Summer 2022: The Same Attack, Two Outcomes]
The clearest demonstration of why phishing-resistant authentication matters is a controlled
comparison rather than a single failure. In summer 2022 a campaign later named 0ktapus by
Group-IB researchers sent near-identical SMS phishing messages to employees at both Twilio and
Cloudflare, each linking to a convincing replica of the companies' single sign-on login
page.\srccite{1.2}{3} Twilio was compromised: employees entered credentials and live one-time
codes, and the attacker relayed both in real time to reach internal systems. Cloudflare was not.
At Cloudflare, three of 76 employees who received the message did enter their username and
password into the fake page -- the human layer failed at almost exactly the rate the previous
section's misconceptions warn about. But Cloudflare does not issue TOTP codes to its employees at
all; every employee authenticates with a FIDO2 hardware security key instead. A key issued for the
real single sign-on domain simply produces no cryptographic assertion for a look-alike domain, so
there was no live one-time code for the phishing page to capture in the first place, and no
dialogue box to fool.\srccite{1.2}{4} The compromise of three sets of credentials produced zero
account takeovers, because origin binding turns a successful phish into an authentication that
still cannot complete. This is the strongest argument in this chapter for treating multi-factor
authentication as a spectrum rather than a checkbox. A TOTP code and a hardware security key both
satisfy ``something you have,'' but only one of them is bound to the domain requesting it -- and
only one of them can be phished out of a human at all.
\end{example}

\section*{Review Questions}
\begin{reviewquestions}
\item An authentication log shows six failed logins for user \texttt{jsmith} between 03:14:02 and
03:14:47, followed by a success from a device never seen before. Explain which detection signals
this triggers and why the combination is stronger evidence than any one signal alone.
\item Explain, in terms of search-space math, why a 20-character passphrase of ordinary English
words can be stronger than a 10-character password that mixes uppercase, lowercase, digits, and
symbols.
\item A password policy requires at least one uppercase letter, one digit, and one symbol, with a
minimum length of 8. Explain two ways this policy can still permit passwords that a dictionary
attack cracks in seconds.
\item Walk through, step by step, how a TOTP verifier decides whether a submitted 6-digit code is
valid for a given timestamp, and explain why the code becomes permanently invalid after 30
seconds rather than merely ``expiring.''
\item Explain why SMS-delivered one-time codes are considered weaker than TOTP codes even though
both appear, to the user, as a 6-digit number they type into a form.
\item In the Cloudflare vs.\ Twilio comparison, three Cloudflare employees typed their username and
password into the same style of fake login page that fully compromised Twilio, yet no Cloudflare
account was taken over. Identify exactly which property of a FIDO2 hardware key made the
difference, and explain why a TOTP code does not have that property.
\item Colonial Pipeline's breach began with a VPN account that was never decommissioned and never
enrolled in MFA. Explain why this account was more dangerous than an active employee account with
the same weak, reused password, even though both would technically be vulnerable to the same
credential-stuffing attempt.
\end{reviewquestions}

\begin{answerkey}
\item Six failures in under a minute triggers the repeated-failed-attempts signal, and a login
succeeding from a device with no prior history triggers the unknown-device signal; the 03:14 hour
also plausibly triggers an unusual-time deviation from a normal baseline. Any one signal alone has
an innocent explanation -- a typo, a new phone, a night owl -- but the combination, arriving
together and culminating in a success, is a pattern that correlation-based detection treats as
strong evidence precisely because independent weak signals reinforcing one another are unlikely to
occur together by chance.
\answerlabel{1.2.A}
\item A 20-character passphrase drawn from ordinary words still multiplies its search space by a
full factor for every added character, and twenty characters is a much larger exponent than ten,
even though each individual word is drawn from a comparatively small dictionary. A 10-character
mixed-character password gains only a modest, one-time widening of the alphabet per position; it
does not compound the way additional length does. Length dominates the calculation because the
exponent, not the base, is what the search space scales on.
\answerlabel{1.2.B}
\item First, the policy sets no restriction on dictionary membership, so a common word satisfying
every rule (for example \texttt{Password1!}) is one of the first candidates any mangling engine
tries. Second, the eight-character floor is short enough that even a fully random string drawn from
the required character classes sits within a search space crackable in a practical amount of time
on modern hardware, especially if the underlying hash is unsalted or uses a fast, non-memory-hard
algorithm.
\answerlabel{1.2.B}
\item The verifier computes the current 30-second time step from the submitted timestamp, computes
\texttt{HMAC-SHA1(secret, time\_step)}, applies dynamic truncation to obtain a 6-digit numeric
code, and compares that value to the code the user submitted using constant-time comparison. The
code becomes permanently invalid, not merely stale, because the next time step produces an
entirely different HMAC output; there is no relationship between consecutive codes that would let
an old code be extended or reused, so ``expired'' is really ``computed for an interval that has
already passed and will never recur.''
\answerlabel{1.2.C}
\item An SMS code travels over the carrier network and can be intercepted, redirected through SIM
swapping, or captured on a compromised carrier system -- weaknesses that exist independent of
anything the user does correctly. A TOTP code requires no network transmission at verification
time at all: both sides compute it locally from a shared secret established once at enrollment, so
there is no carrier-network attack surface to exploit, even though a phished TOTP code can still be
relayed by an attacker fast enough to beat its 30-second window.
\answerlabel{1.2.C}
\item A FIDO2 hardware key cryptographically binds its response to the exact origin (domain)
requesting it, so a look-alike domain receives no valid assertion at all -- there is nothing for
the key to hand over, no matter how convincing the page looks. A TOTP code has no concept of
origin: it is just a number derived from time and a shared secret, so a human who is fooled into
typing it into the wrong page hands the attacker a code that works anywhere, for a few seconds,
including on the real site the attacker is simultaneously relaying it to.
\answerlabel{1.2.C}
\item An active employee account is watched: unusual logins, failed attempts, and behavioral
deviations all have an established baseline to compare against, and a security team reviewing
active accounts would eventually notice something amiss. A dormant, never-decommissioned account
has no active-user review to catch it and no behavioral baseline to deviate from, so a successful
login on it triggers no anomaly at all -- the same weak password is more dangerous there because
nothing in the monitoring pipeline is positioned to notice it being used.
\answerlabel{1.2.A}
\end{answerkey}

\section*{Sources}
\begin{enumerate}[label=\lbrack\arabic*\rbrack, leftmargin=2.2em]
\item \srcitem{1.2}{1}Joseph Blount (Colonial Pipeline CEO), testimony before the U.S.\ Senate
Committee on Homeland Security \& Governmental Affairs, June 8, 2021.
\item \srcitem{1.2}{2}U.S.\ Department of Justice, ``Department of Justice Seizes \$2.3 Million in
Cryptocurrency Paid to the Ransomware Extortionists Darkside,'' press release, June 7, 2021.
\item \srcitem{1.2}{3}Group-IB, ``Roasting 0ktapus: The Phishing Campaign Going After Okta Identity
Credentials,'' Group-IB Blog, August 25, 2022.
\item \srcitem{1.2}{4}Cloudflare Blog, ``The Mechanics of a Sophisticated Phishing Scam and How We
Stopped It,'' July 2022.
\end{enumerate}
